\documentclass[11pt]{article}
\usepackage{amsmath,amssymb}
\usepackage[hidelinks]{hyperref}
\title{Commit the report behind the work counted at the cut}
\author{}
\date{}
\begin{document}
\maketitle

\paragraph{Claim and scope.}
The old, mandatory report is the exact pinned MRWS completion carried in a
173-byte packet. A designated SubT supplier \(Q\) may have a 60-byte optional
probe queued or charged already. A ready \(R\) is preauthorized as fallback.
The appointment, rejoin, venue paths, single Base receiver and optional
withdrawal permission are authored; actual sender debits and callbacks come
from the selected broker/RF code. No field delivery or stock API is asserted.

\paragraph{Interface.}
Preparation can copy or inspect evidence without stopping dispatch.
The consuming adapter takes the same mutex as native queue dispatch, scans
the actual ledger and extant queue, removes the exact pending optional
packet only if authorized and useful, and appends the chosen mandatory
send in that same critical section. It does not wait for a pose, run the
dispatcher, or call a client while holding this added critical section.
An observed receipt, a prior snapshot, or a successful queue removal
alone does not replace the sender's debit. Negative controls use ordinary
scripted dispatch orders in the deterministic fixture.

\paragraph{Prefix certificate.}
Let \(t\) be the append cut, \(L\) the live \(Q\)-charge counted at \(t\),
and \(b_i m_i\) the length times number of candidate receiver calls for
earlier entry \(i\) from \(Q\), after any authorized removal. The guard is
\[
  L+\sum_{i\text{ ahead from }Q} b_i m_i + 173 \;\le C .
\]
The native effect is an attempt per bound recipient, not one per queue entry.
Every later enqueue is behind the report. For its actual attempt at or after
\(t\), remaining pre-cut charge is at most \(L\). All intervening charge ahead
of it is bounded by the retained prefix sum. These two facts prove the
sender allowance, without keeping the lock from append until actual dispatch.
This uses FIFO, nonexpanding endpoint/attempt sets, a monotone clock and no
other charge source. In the selected positive envelope there is also a working
route and a separately uncongested Base with a dispatch by the deadline.
A sufficient bound on sender attempts is not itself evidence of reception.
The sum need not be necessary if some predecessor will be skipped, or expire
before it can interfere. It is sharp for a retained prefix that can all
charge while the old \(L\) remains live.

\paragraph{Discriminating cases.}
A normal dispatch between preparation and publication can move the original
probe from pending to spent even with no callback. If its charge is still
live, trying to remove it cannot recover the allowance. The guarded append
uses \(R\) when it must. Conversely, a pending 60-byte probe with two bound
receivers can charge 120 bytes to \(Q\); the guard for that retained case
is 293 bytes, not 233. At 292 the two-receiver case requires authorized
removal or the fallback. At 293 the retained queue suffices. Appending
another optional packet behind the report does not change that bound.
An endpoint set that can expand before dispatch, or an insertion ahead
of the report, breaks the prefix argument even with an atomic append.
One would have to cap attempted effects as well or requalify at them.

\paragraph{Comparison.}
At capacities at least 173, 360 traces per arm: optimistic preparation
delivers 272, a coherent guard charging one call per queue entry delivers
336, and the prefix-aware guard delivers 360. Fallback counts are 72, 120,
and 132, respectively; optional removals are 24, 24, and 36. An
always-ready fallback can preserve all tested old reports at the price of
using it for all of them. The same-information direct formula agrees on all
1,296 traces, and disjoint traffic is delivered in all traces. Lean checks
the migration arithmetic, generic prefix implication and the double-receiver
threshold. Scope: selected deterministic interleavings, not a threaded stress
or deployment run. The transaction uses the standard
\href{https://www.cs.columbia.edu/~wing/publications/HerlihyWing90.pdf}{linearizability}
pattern; the useful result is what this particular effect owner must count.

\paragraph{What is charged to repair.}
Approach reads all \(q\) queue rows and \(b\) native \(Q\) debit rows when
evaluating \(Q\), plus per-key endpoint lookup and at most one queue erase.
The replay checks those read counts. For bounded message lengths this is
\(O(q\log(1+e)+b)\) with \(e\) endpoint keys (plus linear worst-case erase).
It holds the existing broker-wide mutex during that work, not a lock just
on \(Q\); the simulation places independent work after publication rather
than proving any wall-clock bound or progress during the native all-pose
stall. Index maintenance under churn, an actual producer for the
authorization, and the independent test remain open. No novelty, I1 or
R2 acceptance is claimed.
\end{document}
